\section{Evaluation Protocol}
\label{sec:evaluation}

\subsection{Conventional metrics}
\label{sec:conventional-metrics}

Following~\cite{singla2022explaining}, the evaluation set is split by classifier confidence into $X_n = \{x : f(x) < 0.2\}$ and $X_p = \{x : f(x) > 0.8\}$, up to $n_{\text{eval}} = 200$ each, and each directional policy generates from the corresponding group with its deterministic (mean) action. Full definitions of every metric below are given in Supp.~Sec.~R6.

\paragraph{Counterfactual Validity.} With $s = 2c - 1$ ($+1$ for n$\to$p, $-1$ for p$\to$n),
\begin{equation}
\mathrm{CV} = \frac{1}{N}\sum_{i=1}^{N} \mathbf{1}\!\left[\, s\,\Delta_i \geq \tau \,\wedge\, B_i \,\right],
\label{eq:cv}
\end{equation}
where $\Delta_i = f(x^i_c) - f(x^i)$, $B_i$ requires $s\,(f(x^i_c) - 0.5) > 0$ and $s\,(f(x^i) - 0.5) \leq 0$, $f$ is the guidance classifier (the ensemble mean in ensemble arms) and $\tau = 0.8$. On $X_n$ and $X_p$ the crossing condition $B_i$ is redundant: $f(x) < 0.2$ and $s\,\Delta_i \geq 0.8$ already give $f(x_c) \geq 0.8$ for n$\to$p, and symmetrically for p$\to$n. The tightening over the original, direction-agnostic definition therefore comes from the sign requirement alone: an image moving the wrong way by a large margin is no longer credited.

\paragraph{FID.} Between each direction's counterfactuals and the real images of its target group ($X_p$ for n$\to$p, $X_n$ for p$\to$n; selected by $f$, not by label), on 1024-d features of a TorchXRayVision DenseNet-121~\cite{cohen2022torchxrayvision} with \texttt{rsna} weights rather than Inception-v3. These weights also guide generation in the single-classifier and 2-model arms, so FID is not independent of guidance.
% TODO: 3-model ensemble membership not recorded

\paragraph{Identity Preservation (IP).} As a dependency-free substitute for FOP we report MS-SSIM and HF-corr, the Pearson correlation of high-pass-filtered images rescaled to $[0, 1]$ (0.5 means uncorrelated), a proxy for whether thin structures such as tubes, leads and pacemakers survive editing.

\paragraph{Edit concentration.} The share of $\sum |x_c - x|$ held by the top 5\% of pixels, averaged over pairs; spatially uniform change scores exactly 0.05. We also report the fraction of pixels with $|x_c - x| > 0.01$. There is no calibrated threshold; as an illustrative heuristic only, a compact lesion-shaped edit would place roughly half or more of its mass in the top 5\%.

\subsection{Ground-truth-anchored Lesion Hit Rate}
\label{sec:lhr}

% TODO: LHR implementation not available; definition below is unverified (per-image vs pooled averaging, multi-box handling, test used for the p-value).
CV, FID and IP share a blind spot: none of them can tell whether the edit landed on the pathology. RSNA ships 9{,}555 radiologist-drawn opacity boxes, which we retain. LHR measures the fraction of an edit's absolute-difference mass falling inside the annotation, against the box's own area as the null; lift (hit rate over null), not raw hit rate, is therefore the informative quantity. LHR is defined only for p$\to$n, since n$\to$p sources are classifier-negative and carry no annotation. Details that the definition leaves open are listed in Supp.~Sec.~R6.

We calibrate the metric before use on synthetic edits with known ground truth: a uniform wash scores lift 1.00 ($p=0.59$, chance, as it should), an on-lesion edit scores 12.4 ($p = 3\times10^{-49}$), an off-lesion edit scores 0.

\subsection{Held-out transfer}
\label{sec:transfer}

Because CV is computed with the guidance classifier, it cannot separate a class-relevant edit from an adversarial perturbation. We therefore re-score every counterfactual with classifiers that never participated in guidance (\texttt{nih}, \texttt{pc}, \texttt{mimic\_ch}) and report, per validator $v$ and direction, the ratio of the validator's mean signed score shift to the guidance model's own shift over the same $N$ images,
\begin{equation}
\rho_v = \frac{\sum_i s\,\big(f_v(x^i_c) - f_v(x^i)\big)}{\sum_i s\,\big(f(x^i_c) - f(x^i)\big)},
\label{eq:transfer}
\end{equation}
a ratio of means rather than a mean of per-image ratios, undefined when the guidance shift is zero. The denominator must be the shift of the classifier(s) that actually guided generation.
% TODO: code not preserved (denominator stamp mechanism)
A ratio below 0.15 we treat as the adversarial regime; this is a heuristic operating point, not a calibrated boundary (\S\ref{sec:cv-overstates}). A ratio of 0.5, meaning validators reproduce half the guidance shift, is used only as an illustration and not as a threshold. % TODO: code not preserved -- shard-stamp denominator. Recorded code takes the guidance roster from cache/--guidance-weights (heldout_validation.py:206-228,426-431); current_cfg_stamp does not record USE_GUIDANCE_ENSEMBLE (metrics.py:141-174).

\subsection{Controls the comparison requires}
\label{sec:controls}

A control-arm comparison needs three preconditions that we had assumed and that did not hold (full account in Supp.~Sec.~E1): (i) one deterministic, persisted, patient-level split shared by LoRA fine-tuning, the PPO pool and evaluation, replacing a shared dataset object under which evaluation images had been seen during fine-tuning; (ii) seeded training, since two identical unseeded submissions produced two different policies; and (iii) control arms exercised end to end, which surfaced a routing bug and a shared output directory that would each have produced a well-formed but wrong table.
% TODO: fix code not preserved (seeding)
% TODO: restore claim if code recovered
Two caveats remain. We have no record that the shared generator was fine-tuned again after the split was introduced, so its fine-tuning set may overlap the evaluation images in every arm.
% TODO: verify whether the generator was re-fine-tuned on the split (recorded code fine-tunes on the full dataset)
And the single-classifier arm of \S\ref{sec:five-arms} predates (i) and (ii), was re-scored with the corrected evaluation code, and probably differs from the ensemble arms in reward weights, $a_1$ range and anchoring as well as in guidance (unverified; its training configuration was not recorded), so it is not evaluated under the same protocol as the other four arms.
